\setupencoding[default=ec] % needed for fourier

% \setupfontsynonym[Serif][handling=quality]
% \setupalign[hanging]

\setuplayout
  [width=middle,
   height=middle,
   header=2cm,
   footer=2cm,
   topspace=1.5cm,
   backspace=2cm]

\usetypescript
  [fourier][ec]

\setupbodyfont
  [fourier,12pt]

\setupcolors
  [state=start]

\setupwhitespace
  [big]

\definelayer[page][width=\paperwidth,height=\paperheight]

\definecolor[darkblue]  [c=1,m=1,y=.3]
\definecolor[darkred]   [c=.3,m=1,y=0]
\definecolor[darkyellow][c=0,m=0,y=.7]

\dostepwiserecurse091{\definecolor[xtag:\recurselevel][black]}

\setupfooter
  [style=\bfa,
   color=darkblue]

\setupheader
  [style=\bfa,
   color=darkblue]

\setupheadertexts
  [pagenumber]

\setupfootertexts
  [XML DIR]

\starttext

\startTEXpage

  \setlayerframed
    [page]
    [preset=middle]
    [frame=off,
     width=\paperwidth,
     height=\paperheight,
     backgroundoffset=10pt,
     background=color,
     backgroundcolor=darkyellow]
    {}

  \setlayer
    [page]
    [preset=righttop,offset=.02\paperheight]
    {\rotate[rotation=270]{\scale[height=.25\paperheight]{\darkblue\bf X{\kern.075em}M{\kern0em}L}}}

  \setlayer
    [page]
    [preset=leftbottom,offset=.02\paperheight]
    {\rotate[rotation=270]{\scale[height=.25\paperheight]{\darkblue\bf D{\kern-0.02em}I{\kern-0.005em}R}}}

  \setlayer
    [page]
    [preset=middle]
    {\rotate[rotation=270]{\scale[height=.05\paperheight]{\darkred\bf Hans Hagen -- Pragma ADE}}}

  \tightlayer[page]

\stopTEXpage

This document discusses a (simple) interface from \CONTEXT\ to the
underlying file system. For the moment this discription is limited to
simple  calls to utilities, the schema, and the modules. A directory
listing is generated as follows:

\starttyping
xmltools --dir --pattern=x*.tex --output=xmldir.xml
\stoptyping

You can do that before/between runs, but also from within \CONTEXT:

\startbuffer
\executesystemcommand
  {xmltools --dir --pattern=x*.tex --output=xmldir.xml}
\stopbuffer

\typebuffer

or when that fails:

\startbuffer
\executesystemcommand
  {texmfstart xmltools --dir --pat=x*.tex --out=xmldir.xml}
\stopbuffer

\typebuffer

or more efficient, only in the first run:

\startbuffer
\startmode[*first]
  \executesystemcommand
    {xmltools --dir --pattern=x*.tex --output=xmldir-1.xml}
\stopmode
\stopbuffer

\typebuffer \getbuffer

This produces a file that (on my system) looks as follows:

\typefile{xmldir-1.xml}

Such a file conforms to the following (relax-ng) specification:

\showRNGcomponent [xmldir.rng] [grammar]

The parent element is defined as follows:

\showRNGcomponent [xmldir.rng] [files]

This element has zero or more children:

\showRNGcomponent [xmldir.rng] [directory]

A directory describes files:

\showRNGcomponent [xmldir.rng] [file]

A file element provides the following information:

\showRNGcomponent [xmldir.rng] [descriptors]

As you can see, we kept things simple. You can influence the data
set with the following options to \type {textools}:

\starttabulate[|Tlf{\tt}|p|]
\NC recurse   \NC recurse into subdirectories \NC \NR
\NC pattern   \NC the files to list           \NC \NR
\NC output    \NC the output file             \NC \NR
\NC url       \NC an url to add to the file   \NC \NR
\NC stripname \NC strip basename from pattern \NC \NR
\NC root      \NC the path to run on          \NC \NR
\stoptabulate

Juts for fun, we will now process that file:

\startmode[*first]
  \executesystemcommand
    {xmltools --dir --pattern=x*.tex --output=xmldir-1.xml}
  \executesystemcommand
    {texexec --use=dir-01 --pdf xmldir-1.xml}
\stopmode

\starttyping
texexec --use=dir-01 --pdf xmldir-1.xml
\stoptyping

This produces \in {figure} [fig:xmldir-1]. As you may expect, the
entries are hyperlinks to the files. This also explains the \type
{--url} option to  \type {xmltools} and the \type {url} attribute. We
use this module to provide controlled access to files (graphics) on
a server.

\placefigure
  [here]
  [fig:xmldir-1]
  {The output of module \type {x-dir-01}.}
  {\externalfigure[xmldir-1.pdf][width=\textwidth]}

\startmode[*first]
  \executesystemcommand
    {xmltools --dir --pattern=../allkind/m*.pdf --output=xmldir-2.xml}
  \executesystemcommand
    {texexec --use=dir-02 --pdf xmldir-2.xml}
\stopmode

\starttyping
texexec --use=dir-02 --pdf xmldir-2.xml
\stoptyping

\startbuffer
\startcombination[2*2]
  {\externalfigure[xmldir-2.pdf][background=color,backgroundcolor=gray,width=.45\textwidth,page=1]} {}
  {\externalfigure[xmldir-2.pdf][background=color,backgroundcolor=gray,width=.45\textwidth,page=2]} {}
  {\externalfigure[xmldir-2.pdf][background=color,backgroundcolor=gray,width=.45\textwidth,page=3]} {}
  {\externalfigure[xmldir-2.pdf][background=color,backgroundcolor=gray,width=.45\textwidth,page=4]} {}
\stopcombination
\stopbuffer

\placefigure
  [here]
  [fig:xmldir-2]
  {The output of module \type {x-dir-02}.}
  {\getbuffer}

If you want information about a file you can use the following
module:

\startbuffer
\usemodule[dir-05]
\stopbuffer

\typebuffer \getbuffer

\startbuffer
\getfilestate{xxmldir.tex}

\starttabulate[|||]
\HL
\NC name \NC \getvariable{filestate}{name} \NC \NR
\NC base \NC \getvariable{filestate}{base} \NC \NR
\NC type \NC \getvariable{filestate}{type} \NC \NR
\NC size \NC \getvariable{filestate}{size} \NC \NR
\NC date \NC \getvariable{filestate}{date} \NC \NR
\HL
\stoptabulate
\stopbuffer

\typebuffer

This will produce:

\getbuffer

\stoptext